% ═══════════════════════════════════════════════════════════════
%  PART XVII — APRIL 8, 2026: GAUSSIAN UNIFICATION + SELECTION
% ═══════════════════════════════════════════════════════════════
\clearpage
\part{The Gaussian Unification: All Couplings from $z=11+4i$}

\section{Lock~9: $\alpha^{-1}=137$ Uniquely Selects $q=3$}

For $\GQ(q,q)$, $(k{-}1)^2+\mu^2 = q^4+2q^3+2$.
Setting this equal to~137:
\begin{equation}
\boxed{q^3(q+2) = 135 = 3^3\times 5.}
\end{equation}
This has the \textbf{unique} positive integer solution $q=3$.
Verified computationally for $q\in\{2,3,4,5,7,8,9,11,13\}$:
no other $q$ gives~$137$.

\begin{theorem}[7/7 Observable Test]
Among all $\GQ(q,q)$, only $q=3$ simultaneously matches:
$\alpha^{-1}\approx 137$,
$\sin^2\theta_{12}\approx 0.307$,
$\sin^2\theta_{23}\approx 0.546$,
$\sin^2\theta_{13}\approx 0.022$,
$\alpha_s\approx 0.118$,
$f = 24$,
$E = 240$.
No other $q$ matches even one observable.
\end{theorem}


\section{Lock~10: The Gaussian Coincidence $k{-}1=\Phi_3{-}\lambda$}

For general $\GQ(q,q)$: $k{-}1=q^2{+}q{-}1$ and $\Phi_3{-}\lambda=q^2{+}2$.
These are equal if and only if $q=3$.

\begin{theorem}[Gaussian Unification]
\label{thm:gaussian}
At $q=3$, the Gaussian integer $z = (k{-}1)+\mu i = (\Phi_3{-}\lambda)+\mu i
= 11+4i$ encodes all three gauge couplings:
\begin{align}
\alpha_{\mathrm{em}}^{-1} &= |z|^2 = 137 \quad (\text{tree level}), \\
\alpha_s &= \mu^2/|z|^2 = 16/137 \quad (\text{tree level, $1.3\sigma$}), \\
\sin^2\theta_W &= q/\Phi_3 = q/((k{-}1)+\lambda) \quad (\text{dressed}).
\end{align}
The coupling ratios at tree level satisfy $\alpha_s/\alpha_{\mathrm{em}} = \mu^2 = s^2$
and
\begin{equation}
\alpha^{-1} = \left(\frac{q}{\sin^2\theta_W} - \lambda\right)^2 + \mu^2.
\end{equation}
This algebraic relation between $\alpha^{-1}$ and $\sin^2\theta_W$
holds \textbf{only at $q=3$}.
\end{theorem}


\section{The 7/7 Definitive Selection Table}

\begin{center}
\begin{tabular}{@{}crrrrrrrr@{}}
\toprule
$q$ & $\alpha^{-1}_0$ & $\sin^2\theta_{12}$ & $\sin^2\theta_{23}$
& $\sin^2\theta_{13}$ & $\alpha_s$ & $f$ & $E$ & score \\
\midrule
$2$ & $34$ & $0.429$ & $0.429$ & $0.059$ & $0.204$ & $9$ & $45$ & $0/7$ \\
$\mathbf{3}$ & $\mathbf{137}$ & $\mathbf{0.308}$ & $\mathbf{0.538}$
& $\mathbf{0.022}$ & $\mathbf{0.118}$ & $\mathbf{24}$ & $\mathbf{240}$
& $\mathbf{7/7}$ \\
$4$ & $386$ & $0.238$ & $0.619$ & $0.009$ & $0.077$ & $50$ & $850$ & $0/7$ \\
$5$ & $877$ & $0.194$ & $0.677$ & $0.004$ & $0.054$ & $90$ & $2340$ & $0/7$ \\
$7$ & $3089$ & $0.140$ & $0.754$ & $0.000$ & $0.031$ & $224$ & $11200$ & $0/7$ \\
\bottomrule
\end{tabular}
\end{center}
\emph{Seven independent observables from seven experiments.
Only $q=3$ matches all seven. No other $q$ matches even one.}
